\chapter{1974 Nobel Prize in Economics}
\textbf{Laureates: }Gunnar Myrdal (Sweden), Friedrich Hayek (Austria)

\section*{Reason for Award}
Institutional analysis price income business cycle fluctuations

\section{What They Did}
Gunnar Myrdal focused circular causation cumulative causation development initial advantages reinforce themselves over time Friedrich Hayek advanced spontaneous order theory arguing complex social orders emerge without deliberate design knowledge utilization decentralized markets critiques central planning based impossibility centralized calculation

\section{Impact on Economic Thinking}
Their contrasting perspectives shaped institutional economics macroeconomic thought regarding market mechanisms planned systems influencing debates government intervention versus market freedom throughout twentieth century Myrdal s work emphasized cumulative process development disparities attention distributional consequences policies Hayeks contribution emphasized price signals coordination function knowledge dispersed society warning central planners ignorance essential information

\section{Modern Perspectives Today}
Myrdal's insights influence contemporary regional development studies wealth disparity analysis Hayeks work on knowledge institutions continues influencing discussions economic freedom institutional evolution limits central planning In transition economies his ideas informed privatization governance reforms Both thinkers continue reference point debates globalization market integration government role market economies
